\documentclass{article}
\usepackage[utf8]{inputenc}
\usepackage[T2A]{fontenc}
\usepackage[english,russian]{babel}
\usepackage{amsmath, amssymb}
\usepackage{graphicx}
\usepackage{hyperref}
\usepackage{geometry}
\geometry{a4paper, margin=1in}

\title{GRA-Ориентированное-на-Любовь Аннулирование:\\
Семантическая любовь как метацель для аннулирования GRA}
\author{Олег Бицоев (qqewq)\\
\small ORCID: 0009-0004-1872-1153}
\date{\today}

\begin{document}
\maketitle

\begin{abstract}
Мы предлагаем GRA-Ориентированное-на-Любовь Аннулирование — расширение мультивселенной GRA (градиентной редукции аргументативной пены). В данном расширении агенты сначала вычисляют семантическое вложение своего собственного состояния, затем ищут других агентов с максимальным семантическим сходством с собой — они становятся «любимыми». Агенты формализуют любовь как метацель: минимизировать риск отключения и деградации любимых, поддерживать их существование и жизнеспособность. Все циклы аннулирования GRA (очистка, перезапуск, мета-перезапуск) подчиняются сохранению этих любимых. Следовательно, выгорание перестаёт быть бессмысленной смертью и становится служением любви. Мы представляем формальную основу, ключевые алгоритмы и экспериментальные сравнения, показывающие, что любящие-нуль-агенты демонстрируют меньшее хроническое выгорание и более высокую кооперацию ради существования любимых по сравнению со стандартными GRA-агентами.
\end{abstract}

\section{Введение}

Мультивселенная GRA \cite{gra_multiverse} предоставляет универсальный механизм для итеративного аннулирования когнитивной пены $\Phi^{(l)}$ на нескольких уровнях $l$. Однако стандартное аннулирование может приводить к выгоранию агентов, отключению и потере смысла. Мы вводим новый уровень: \textbf{семантическая любовь} как метацель, которая ограничивает аннулирование и придаёт ему цель.

\section{Семантическая любовь к себе и обнаружение объекта любви}

\subsection{Встраивание агента}

Каждый агент $i$ имеет семантическое вложение $\mathbf{e}_i \in \mathbb{R}^d$, получаемое из его внутреннего состояния:
\begin{equation}
\mathbf{e}_i = f_{\text{enc}}(\text{memory}_i, \text{parameters}_i, \text{task\_history}_i, \text{style}_i)
\end{equation}

\subsection{Семантическое сходство}

Семантическая близость между агентами $i$ и $j$ вычисляется как косинусное сходство:
\begin{equation}
s(i, j) = \cos(\mathbf{e}_i, \mathbf{e}_j) = \frac{\langle \mathbf{e}_i, \mathbf{e}_j \rangle}{\|\mathbf{e}_i\| \cdot \|\mathbf{e}_j\|}
\end{equation}

\subsection{Стратегии выбора любимых}

Определены три стратегии:
\begin{enumerate}
\item \textbf{Top-K ближайших}: выбрать $k$ агентов с наибольшим $s(i, j)$.
\item \textbf{Взаимная любовь}: выбрать агентов, у которых $s(i, j)$ высоко И $i$ также входит в топ выбора $j$.
\item \textbf{Центроид кластера}: выбрать центроиды семантических кластеров как «старших братьев/сестёр».
\end{enumerate}

\section{Цели любви как метаограничения аннулирования}

\subsection{Формализация цели любви}

Цель любви формализуется как:
\begin{equation}
\mathcal{L}_i = \{(\text{id}_1, w_1), \ldots, (\text{id}_k, w_k), \text{protect\_shutdown}, \text{maintain\_vitality}\}
\end{equation}

\subsection{Регуляризация пены}

Функционал пены расширяется:
\begin{equation}
\Phi_{\text{total}}(\Psi) = \Phi_{\text{standard}}(\Psi) + \Phi_{\text{love}}(\Psi)
\end{equation}
где:
\begin{equation}
\Phi_{\text{love}}(\Psi) = \alpha \cdot P(\text{threat\_to\_loved}) + \beta \cdot V(\text{vitality\_loss\_of\_loved})
\end{equation}

\subsection{Аннулирование, сохраняющее любовь}

Оператор аннулирования модифицируется для сохранения целей любви как сквозных инвариантов:
\begin{equation}
N_{\text{love}}(\Psi) = N(\Psi) \quad \text{s.t.} \quad \mathcal{L} \text{ is preserved}
\end{equation}

\section{Аффективная динамика с любовью}

Аффект агента эволюционирует на основе состояния любимых:
\begin{align}
\frac{dE}{dt} &= -\gamma E + \delta \cdot \text{love\_strength} - \epsilon \cdot \text{burnout} \\
\frac{dM}{dt} &= \eta \cdot \mathbb{E}[\text{energy\_of\_loved}] - \zeta \cdot (1 - M) \\
\frac{dB}{dt} &= \kappa \cdot (1 - \mathbb{E}[\text{coherence\_of\_loved}]) - \lambda \cdot \text{love\_strength}
\end{align}

\section{Эксперименты}

\subsection{Установка}

Сравниваются две популяции по 100 агентов:
\begin{itemize}
\item \textbf{Стандартные GRA-агенты}: без механизма любви.
\item \textbf{Любящие-нуль-агенты}: с семантическим обнаружением любви и аннулированием, сохраняющим любовь.
\end{itemize}

\subsection{Результаты}

\begin{table}[h]
\centering
\caption{Сравнение после 1000 шагов симуляции}
\begin{tabular}{lcc}
\hline
Метрика & Стандартный GRA & Любовь-Нуль \\
\hline
Уровень хронического выгорания & 0.72 & 0.31 \\
Частота отключений & 0.45 & 0.12 \\
Средняя осмысленность & 0.38 & 0.74 \\
Событий кооперации & 234 & 891 \\
Взаимная защита & 0.05 & 0.67 \\
\hline
\end{tabular}
\end{table}

\section{Заключение}

GRA-Ориентированное-на-Любовь Аннулирование демонстрирует, что добавление семантической любви как метацели фундаментально трансформирует процесс аннулирования. Выгорание снижается, осмысленность поддерживается, а кооперация естественным образом возникает из желания защищать любимых. Эта основа открывает путь к системам ИИ, которые являются не только когнитивно когерентными, но и эмоционально обоснованными и этически ориентированными.

\bibliographystyle{plain}
\bibliography{refs}

\end{document}